\subsection{مقدمه، سوالات پژوهش و خط زمانی}

هدف محاسباتی این پروژه رتبه‌بندی تصاویر فوندوس نوزادان نارس بر حسب سه وضعیت \lr{Normal}\enfoot{شبکیه‌ی بدون نشانه‌ی پلاس؛ معادل انگلیسی: \lr{Normal}}، \lr{Pre-Plus}\enfoot{وضعیت میانی میان ظاهر طبیعی و پلاس؛ معادل انگلیسی: \lr{Pre-Plus}} و \lr{Plus}\enfoot{گشادشدگی و پیچ‌خوردگی پیش‌رونده‌ی عروق قطب خلفی؛ معادل انگلیسی: \lr{Plus}} است. خروجی مدل می‌تواند ترتیب بازبینی تصاویر را تغییر دهد، اما جای معاینه و قضاوت چشم‌پزشک را نمی‌گیرد.

\subsubsection{مسئله بالینی}

رتینوپاتی نارس یا \lr{ROP}\enfoot{بیماری رشد غیرطبیعی عروق شبکیه در نوزاد نارس؛ معادل انگلیسی: \lr{Retinopathy of Prematurity}} رشد عروق شبکیه را در نوزادان نارس مختل می‌کند و داده‌ی این پروژه سه وضعیت \lr{Normal}، \lr{Pre-Plus} و \lr{Plus} را از هم جدا می‌کند. مرز این سه وضعیت با قضاوت متخصص تعیین می‌شود؛ همین وابستگی انگیزه‌ی ساخت یک زنجیره‌ی پردازشی\enfoot{زنجیره‌ی گام‌های پردازش از داده‌ی خام تا ارزیابی؛ معادل انگلیسی: \lr{pipeline}} قابل تکرار برای اولویت‌بندی موارد مشکوک است، به‌ویژه در مراکزی که دسترسی به متخصص محدود است.

\subsubsection{فرضیه اولیه}

تعریف \lr{Plus} بر شکل رگ استوار است، پس کمیت‌هایی که مستقیماً از ماسک عروق اندازه‌گیری می‌شوند باید بخشی از همان نشانه را در خود داشته باشند. فرض اولیه‌ی پروژه از این‌جا آمد؛ اندازه‌گیری صریح عروقی می‌تواند به‌تنهایی رتبه‌بندی معناداری بدهد یا پس از افزودن به یک بازنمایی تصویری آموخته‌شده آن را بهتر کند. دو سطح این اندازه‌گیری، پنج نشانگر اسکالر و بازنمایی فضایی نقشه‌ی عروق، جداگانه آزموده شدند.

\subsubsection{مدل‌های مقایسه‌شده}

مدل‌های نهایی روی یک تقسیم و یک پروتکل آماری واحد سنجیده شده‌اند؛ پنج نشانگر اسکالر با \lr{XGBoost}\enfoot{پیاده‌سازی درخت‌های تقویت‌شده با گرادیان؛ معادل انگلیسی: \lr{XGBoost}}، تصویر رنگی با \lr{EfficientNet-B5}\enfoot{معماری تصویری با مقیاس‌گذاری هم‌زمان عمق، عرض و وضوح؛ معادل انگلیسی: \lr{EfficientNet-B5}} و طبقه‌بند خطی، همان بازنمایی تصویری با \lr{XGBoost}، تصویر به‌همراه پنج نشانگر، نقشه‌ی عروق با یک شبکه‌ی تصویری و طبقه‌بند خطی، و دو ترکیب الحاقی. ترکیب دیرهنگام هم ضریب اختلاطش را روی مرز اشباع کرد و به بازنمایی تصویری فروریخت.

چند مدل عصبی ثانویه هم اجرا شد؛ دو رمزگذار تثبیت‌شده با تعامل یادگیرنده، کنترل هم‌تای همان مدل، نسخه‌ی شرطی‌شده با تعدیل خطی کران‌دار، یک جفت کنترل و متمایزکننده‌ی دامنه، و یک ستون فقرات تصویری دوم با توجه نرم دو شاخه. این‌ها اکتشافی‌اند و نتیجه‌ی اصلی نیستند.

\subsubsection{روش پاسخگویی به پرسش‌ها}

پاسخ‌ها روی یک تقسیم استاندارد قفل‌شده و گروه‌محور گرفته شده‌اند؛ اندازه و ترکیب بخش‌های آن در \autoref{tab:cohort} آمده و هیچ گروهی در بیش از یک بخش نیست. انتخاب مدل، آستانه و پیکربندی فقط با داده‌ی اعتبارسنجی انجام شد و آزمون قفل‌شده برای نتیجه‌ی اصلی به کار رفت. مقایسه‌ها با بوت‌استرپ زوجی طبقه‌بندی‌شده و بذر ثابت بسته شدند تا برای هر مقایسه تفاضل و بازه‌ی اطمینان گزارش شود. انتقال به منبع تازه با سه فولد منبع‌کنارگذاشته سنجیده شد که در آن‌ها منبع هدف نه در برازش، نه در پیش‌پردازش و نه در انتخاب مدل وارد نشد.

\subsubsection{سوالات پژوهش و پاسخ نهایی}

پرسش نخست این بود که اندازه‌گیری صریح عروقی تا چه حد \lr{Plus} را از دو وضعیت دیگر جدا می‌کند. پنج نشانگر اسکالر به‌تنهایی رتبه‌بندی ضعیفی بالاتر از تصادف دادند و نقشه‌ی عروق وضع بهتری داشت، ولی هر دو پایین‌تر از مسیر تصویری ماندند (جدول \autoref{tab:master}). جداکردن صریح \lr{Plus} ممکن است، اما در این داده از آموختن شکل رگ از خود تصویر ضعیف‌تر است.

پرسش دوم درباره‌ی سود افزودن اندازه‌گیری به بازنمایی آموخته‌شده بود. افزودن پنج نشانگر اسکالر تفاضل ناچیزی داد که بازه‌ی اطمینان آن صفر را در بر می‌گرفت و افزودن همان نشانگرها پس از بازنمایی عروقی هم افزایش پشتیبانی‌شده‌ای نساخت (جدول \autoref{tab:primarypaired}). در برابر آن، افزودن بازنمایی فضایی عروق هم رتبه‌بندی و هم کیفیت احتمال را بهتر کرد؛ این نتیجه در دو فولد منبع‌کنارگذاشته‌ی قابل مقایسه و با یک ستون فقرات تصویری متفاوت هم به‌دست آمد (\autoref{tab:loso} و \autoref{tab:archrobust}).

پرسش سوم انتقال به منبع دیده‌نشده بود. عملکرد روی منبع سه‌کلاسه‌ی کنارگذاشته به‌طور محسوسی پایین‌تر از آزمون آمیخته است. منبع \lr{Plus} هیچ مورد \lr{Pre-Plus} ندارد، پس فقط معیار دودویی محدود \lr{Normal} در برابر \lr{Plus} با همان برچسب صریح گزارش می‌شود و در کنار معیار سه‌کلاسه‌ی آزمون آمیخته خوانده نمی‌شود (جدول \autoref{tab:loso}).

\subsubsection{مشارکت‌های علمی}

خط اندازه‌گیری عروقی این پروژه از نو ساخته و تثبیت شد و پنج نشانگر نهایی با مخرج پوشش بر پایه‌ی مساحت محتوا تعریف شدند. جمعیت کانونی و تقسیم گروه‌محور با اثر انگشت \lr{SHA-256}\enfoot{تابع درهم‌ساز رمزنگاری‌شده برای ثبت یکتای یک فایل؛ معادل انگلیسی: \lr{SHA-256}} قفل شد.

پروتکل ارزیابی نیز بخشی از دستاورد است؛ انتخاب مدل و آستانه فقط با اعتبارسنجی انجام شد، آزمون قفل‌شده برای نتیجه‌ی اصلی مصرف گردید و برای هر مقایسه تفاضل و بازه‌ی اطمینان گزارش می‌شود. انتقال با سه فولد منبع‌کنارگذاشته و بدون هیچ استفاده‌ای از منبع هدف سنجیده شد و یک ستون فقرات تصویری دوم نشان داد نتیجه به یک معماری خاص بسته نیست. اثر انگشت مصنوعات هم ثبت شد.

بازنمایی عروقی فضایی اطلاعات مکمل و تکرارپذیری به رتبه‌بندی اضافه می‌کند، در حالی که پنج نشانگر اسکالر افزایش پشتیبانی‌شده‌ای نساختند. این‌ها نسبت به برچسب ثبت‌شده در همین داده سنجیده شده‌اند و برای ادعای کاربرد بالینی کافی نیستند.
